\section{Introduction}
Hypersonic flow over a blunt body produces a detached compression layer in which momentum redirection, translational heating, internal-energy relaxation and gas--surface interaction occur over coupled spatial scales. When the molecular mean free path is no longer negligible relative to the body dimension, the layer has finite kinetic thickness and the continuum notion of a discontinuous shock surface becomes operational rather than exact. DSMC remains the reference particle method in this regime \citep{Bird1994,Lofthouse2007}, and the circular cylinder remains a canonical geometry because it contains a curved detached front, a stagnation region and strong wall interaction without geometric complexity. Shock extraction in such fields is itself non-trivial: gradient- and schlieren-based definitions agree when a localized front exists, but become diagnostic rather than unique as the layer becomes diffuse \citep{Akhlaghi2017,Akhlaghi2021}.

Unsteady shock motion has a long history outside rarefied-gas dynamics. \citet{Plotkin1975} represented a shock as a linearly damped displacement driven by broadband upstream fluctuations, yielding a first-order low-pass response. Experiments and simulations of separated shock-wave/boundary-layer interactions subsequently developed this picture and identified both restoring dynamics and low-frequency forcing associated with the separated region \citep{PoggieSmits2001,PoggieSmits2005,Dussauge2006,Dupont2006,Piponniau2009,TouberSandham2011,PriebeMartin2012}. The modern review of \citet{ClemensNarayanaswamy2014} emphasizes that several mechanisms can generate low-frequency shock motion and that a measured spectrum alone does not identify the forcing. These studies concern turbulent or separated continuum flows, not a rarefied bow shock. Their relevance here is therefore structural rather than literal: they establish a fluid-mechanical question that is sharper than mode detection alone---does the shock layer behave as a damped, spatially organized filter, and what can be inferred about its relaxation time and forcing?

The forcing environment is fundamentally different in DSMC. Molecular and simulator-particle fluctuations are intrinsic to particle fields, and their sampling error depends on particle number, averaging time and the observable being estimated \citep{Stefanov2000,Hadjiconstantinou2003}. Fluctuating-hydrodynamic formulations make explicit that thermal forcing can generate temporally and spatially correlated macroscopic fluctuations \citep{BellGarciaWilliams2007,WilliamsBellGarcia2008,Bell2022}. Molecular simulations further show that stochastic microscopic forcing can seed or alter organized hydrodynamic dynamics, including instability growth and symmetry breaking \citep{Gallis2016,Gallis2017,Gallis2021}. Consequently, the presence of stochastic forcing does not make an observed slow response unphysical, but its amplitude cannot be interpreted independently of simulator-particle weight. The physically transferable quantities are instead the response shape, relaxation time and transfer behaviour, provided they survive changes in sampling and numerical realization.

The specific question here is whether the detached bow layer possesses such a collective displacement coordinate. Translation of a thin or moderately diffuse front produces the first-order disturbance
\begin{equation}
 q'(s,\theta,t)\simeq-a(\theta,t)\frac{\partial \mean q}{\partial s},
 \label{eq:translation}
\end{equation}
where $s$ is distance from the wall along a body-normal ray. If the layer moves collectively, $a(\theta,t)$ should contain a slowly varying component shared by distant portions of the bow. In DSMC, however, the observed marker also contains finite-particle noise, and the noise is temporally correlated because particle states and sampling accumulators persist between neighbouring outputs. The challenge is therefore to distinguish a low-pass physical response from correlated feature error rather than to ask whether the instantaneous field is low rank.

Registration and modal analysis introduce additional pitfalls. Registration methods can remove transport-dominated variability and improve reduced-order compactness \citep{RowleyMarsden2000,Reiss2018,Taddei2020}, but their validity depends on the map and its physical support. A shock-attached grid that crosses the solid body can be filled by repeated wall-adjacent values and generate artificial low rank. POD identifies variance-optimal structures \citep{Lumley1967,Sirovich1987,Berkooz1993,Taira2017}; DMD and SPOD target temporal or frequency-dependent organization \citep{Schmid2010,Towne2018,SchmidtColonius2020}. Power spectra estimated by Welch averaging are useful diagnostics \citep{Welch1967}, but none of these decompositions alone separates physical covariance from particle-sampling covariance. The present analysis therefore treats modal tools as diagnostics and bases the physical classification on coarse-graining response, out-of-sample covariance prediction, long-range angular correlation, synthetic controls and full-field reconstruction.

Kinetic studies provide the closest rarefied-flow context. DSMC has revealed low-frequency molecular fluctuations inside one-dimensional shocks and analytical models have reproduced part of their macroscopic structure \citep{Sawant2021,Sawant2022}. Kinetic base states have also been combined with linear stability analysis in hypersonic boundary layers and rarefied shear flows \citep{Klothakis2022,Zou2023}, while time-resolved DSMC and data-driven decompositions have exposed organized dynamics in strongly separated high-speed configurations \citep{Karpuzcu2025}. More recently, \citet{Karpuzcu2026} linearized a Boltzmann--BGK model about one-dimensional normal shocks and showed that non-Maxwellian velocity distributions move spectral branches towards less stable locations. That analysis concerns deterministic perturbations of planar argon shocks, whereas the present work concerns stochastic fluctuations of a curved, wall-coupled nitrogen bow layer. It nevertheless supports a key interpretation: a stable kinetic operator can possess a slowly relaxing macroscopic projection that is continually excited by stochastic forcing.

A separate published study used the same canonical cylinder geometry for neural-operator development, including a steady fixed-Mach-10 Knudsen-number parameterization and wall-response surrogates \citep{RoohiShojaAzghadi2026PoF}. That work established the availability and surrogate learnability of cross-regime steady fields; it did not address temporal shock motion or separate physical covariance from DSMC sampling covariance. The distinction matters for publication overlap: the present contribution is not another surrogate or mean-field study, but an analysis of ordered nitrogen output blocks and their time-dependent collective response.

The companion study of this cylinder family \citep{RoohiShoja2026Mean} examines mean standoff, mean 10--90 thickness, variable-specific relaxation lengths and parameter-indexed shock-attached POD. It establishes that the mean compression layer inflates and loses single-scale similarity as rarefaction increases. None of those parameter-space results establishes that the layer moves collectively in time. The present paper uses consecutive nitrogen outputs at fixed operating conditions and makes four distinct contributions. First, it audits shock-attached support and quantifies false low-rank structures created by solid-side interpolation. Second, it performs temporal coarse graining before marker extraction and separates a persistent covariance from correlated sampling covariance using small-sample model selection, block resampling and synthetic controls \citep{Akaike1974,HurvichTsai1989,Kunsch1989,PolitisRomano1994,Higham2002}. Third, it identifies a common long-range angular displacement shape in the independently analysed $\KnD=0.01$ and $0.025$ records. Fourth, it validates the coordinate in density, Mach number, pressure and translational temperature through the translation template in \cref{eq:translation}. The central claim is deliberately narrow: a weak, low-pass collective response is resolved in two near-continuum states. The data do not establish a disappearance threshold at larger Knudsen number. The two resolved records are used to test consistency, not to claim a universal rarefaction trend.
